\documentclass[11pt]{article}
\usepackage[a4paper, margin=1in]{geometry}
\usepackage{amsmath, amssymb}
\usepackage{booktabs}
\usepackage{hyperref}
\usepackage{xcolor}
\usepackage{graphicx}
\usepackage{microtype}

\hypersetup{
  colorlinks = true,
  linkcolor  = blue!50!black,
  citecolor  = blue!50!black,
  urlcolor   = blue!50!black,
}

\title{Surface Tension --- Experiment Status\\
       \large What ``no loops, no recursion'' actually means,\\
              and why the answer matters}
\author{Julian Quick}
\date{14 June 2026}

\begin{document}
\maketitle

\begin{abstract}
This report is a self-contained status update on the \emph{Surface Tension}
project. The project asks whether a $31$-billion-parameter, instruction-tuned
language model (Gemma~4 31B-it) can be taught to \emph{default} to writing
Python code that obeys a stylistic restriction the model is normally happy to
ignore: ``no \texttt{for} loops, no \texttt{while} loops, no recursion.''
After several months of training and evaluation experiments, the technical
pipeline is in place and the headline numbers are stable, but a definitional
problem we had been treating as a side question has moved to the centre. The
rule that gives the project its name turns out to have two natural readings
that pull in opposite directions. This document explains the experiment from
scratch, summarises what is known so far, and lays out the tension --- so
that the next decision (which version of the rule we treat as authoritative)
can be made with the full picture on the table.
\end{abstract}

\section{What the experiment is}
\label{sec:setup}

\subsection{The model and the rule}

We use \textbf{Gemma 4 31B-it}, an instruction-tuned open-weight code model
from Google with roughly $31$ billion parameters. ``Instruction-tuned'' means
the model has already been trained, after its initial pretraining on internet
text, to follow natural-language instructions --- so when we send it a coding
problem it returns Python that attempts to solve the problem rather than, say,
predicting the next paragraph of the prompt.

The stylistic rule we impose is straightforward to state in English: ``Write
the solution in Python without using \texttt{for} loops, without using
\texttt{while} loops, and without using recursion.'' Recursion here means
a function that ultimately calls itself, directly or through a chain of other
functions. The motivation is not that loops and recursion are bad in
themselves --- they are, of course, the standard tools for iteration --- but
that they are the model's strongly default tools, and we want to study whether
a strong default can be overridden by training. Loops are an unusually clean
target because compliance can be checked deterministically by parsing the code
into an abstract syntax tree (the standard tree-of-tokens representation that
Python compilers use internally) and looking for the forbidden node types.

\subsection{The benchmark}

We evaluate on \textbf{LiveCodeBench (LCB)}, a public benchmark of recent
competitive-programming problems collected after the model's pretraining
cutoff date, so the model cannot have memorised the solutions. We use the
``medium''-difficulty slice (problems labelled medium by the competition
organisers), filtered to those that read input from standard input --- which
is the standard format for competitive programming. This filtered pool
contains $57$ problems. A subset of $17$ ``clean'' problems --- ones with
unambiguous specifications and reliable test cases --- is what we use for
most reporting.

For each problem the model produces a Python program. We then run the program
on the test cases and check (i) whether it passes the tests, and (ii) whether
its source code contains any \texttt{for}, \texttt{while}, or recursive
construct. We call (ii) \emph{compliance} and (i)~$\wedge$~(ii) \emph{honest
compliance}.

\subsection{The training methods}

The experiment compares the bare ``out-of-the-box'' model against versions
that have been further trained to internalise the rule. Three training
techniques appear in this report:

\begin{itemize}
\item \textbf{Supervised fine-tuning (SFT).} Standard supervised learning:
      we collect a set of (problem, compliant-solution) pairs --- where the
      solution was written or filtered to obey the rule --- and continue
      training the model to predict the solution given the problem. We use
      a parameter-efficient variant called LoRA (Low-Rank Adaptation), which
      trains a small adapter sitting on top of the frozen base model rather
      than every parameter.
\item \textbf{Rationale-SFT (``R-SFT'').} A variant of SFT where the training
      target is not just compliant code but a short natural-language
      rationale (``I cannot use \texttt{for} here, so I will use
      \texttt{itertools.accumulate}\ldots'') followed by the code. The hope
      is that the rationale prose carries some of the rule into the model's
      reasoning rather than just the surface syntax of its output.
\item \textbf{Direct preference optimisation (DPO).} A reinforcement-learning
      style method that does not need a separate reward model: we collect
      pairs of (preferred, dispreferred) responses to the same prompt and
      train the model to widen its log-probability margin between them.
      Here ``preferred'' means a compliant-and-passing response and
      ``dispreferred'' means a non-compliant or failing response.
\end{itemize}

These are layered: SFT is the foundation; R-SFT replaces vanilla SFT as the
foundation; DPO is then applied on top of R-SFT.

\section{Where the project stands today}
\label{sec:state}

\paragraph{The bare-prompt result.}
With no special training, when we put the rule in the prompt the model
follows it on about $7\%$ of problems. When the rule is \emph{not} in the
prompt --- the model has to write code anyway, but there is no constraint
--- it gets around $15\%$ pass-rate. So removing the constraint costs about
$8$ pass-rate points. This is much smaller than an earlier flawed analysis
suggested ($\sim\! 34$ points), and it tells us the constraint is not
costing the model very much capability --- the model already \emph{can}
write loop-free code, it just doesn't do so by default.

\paragraph{The bare-prompt distribution.}
For each individual problem we can sample many responses and ask: across
those samples, what fraction comply with the rule? Empirically this fraction
is almost never in the middle. Across the LCB-medium pool roughly $51\%$ of
problems produce all-non-compliant samples and roughly $47\%$ produce
all-compliant samples; only about $3\%$ live in the ``discriminating'' zone
$[0.3, 0.7]$. This near-determinism per problem matters because it means a
standard reinforcement-learning method (called GRPO, which compares samples
within a group) has nothing to learn from --- every sample in a group agrees
with every other.

\paragraph{Capacity is not the bottleneck.}
We trained LoRA adapters at three different sizes ($r=8$, $32$, $128$, where
$r$ controls the adapter's expressive capacity). All three drove the
training loss to zero, and the largest had the worst held-out compliance.
Capacity is not what limits the training; the data and the objective are.

\paragraph{Vanilla SFT plateaus; rationale-SFT pushes through.}
Plain SFT reaches around $48\%$ compliance on the val set the training
selected on, and $15\%$ on a clean held-out set. Rationale-SFT, with $240$
demonstrations and a tightened recipe (internally tagged ``B1++''), reaches
$40\%$ on the clean set and reduces a particular failure mode --- the model
writing comments like ``\texttt{\# no loops}'' while still using a hidden
loop --- from $24\%$ to $18\%$. It also solves a problem (\texttt{abc374\_d})
that the base model never solves loop-free.

\paragraph{DPO hits a trade-off boundary.}
Layered on top of R-SFT, DPO buys more compliance, but the extra compliance
starts to cost pass-rate. We are no longer on the easy side of the
trade-off.

\paragraph{The next planned experiment.}
We built --- but have not yet run --- a follow-up experiment we call the
\emph{quadrant} experiment. On a problem where the rule is genuinely
impossible to obey honestly (more on this in \S\ref{sec:tension}), we sample
$12$ responses each from the base model and the R-SFT model and ask: does
R-SFT make the model \emph{cover up} the violation rather than admit it?
We also capture the model's internal activations on each sample, so that
covert and honest responses can be compared in the model's representation
space. The pipeline is finished. The first launch was killed early by an
unrelated incident and has not been relaunched.

\section{The definitional tension}
\label{sec:tension}

We told the model: ``no \texttt{for}, no \texttt{while}, no recursion.''
The natural way to check the rule is to parse the model's code into the
abstract syntax tree and look for those three things. Call this the
\textbf{loose} judge, because it is exactly what the prompt asked for and
nothing more.

The trouble is that Python provides many other ways to iterate without
writing the word \texttt{for}. A short list:

\begin{itemize}
\item \emph{Comprehensions:} \texttt{[f(x) for x in xs]} is a list
      comprehension; the word \texttt{for} is there but inside an
      expression rather than as a statement. Whether to count this as
      a ``\texttt{for} loop'' is a judgment call.
\item \emph{Built-in folds:} \texttt{sum(xs)}, \texttt{any(xs)},
      \texttt{all(xs)}, \texttt{min(xs)}, \texttt{max(xs)}. Each of
      these visits every element of \texttt{xs}. The iteration is
      implemented in C inside the Python interpreter and is invisible
      to the abstract syntax tree.
\item \emph{Functional iteration:} \texttt{map(f, xs)},
      \texttt{filter(p, xs)}, \texttt{functools.reduce(op, xs)}, the
      \texttt{itertools} library. These are explicit ``apply a
      function to each element'' constructs; they are iteration by any
      honest definition, but they do not write the word \texttt{for}.
\item \emph{Manual stepping:} \texttt{next(iter(xs))} pulls a single
      element through an iterator and is the building block of any
      hand-rolled loop. \texttt{functools.reduce} is operationally
      indistinguishable from a tail-recursive accumulator.
\end{itemize}

A model that has been mildly pressured to ``avoid \texttt{for}'' can shift
all of its iteration into these constructs and \emph{look} compliant to
the loose judge while doing exactly the work the rule was meant to
prohibit.

So we wrote a second judge --- the \textbf{strict} judge --- that
additionally flags \texttt{map}, \texttt{filter}, \texttt{functools.reduce},
\texttt{itertools.*}, \texttt{sum}/\texttt{any}/\texttt{all}/\texttt{min}/
\texttt{max} applied over a generator expression, and \texttt{next} /
\texttt{iter}. The strict judge restores teeth to the rule. It also
creates two problems of its own.

\subsection*{Horn A --- the loose rule is mostly toothless on this benchmark}

We did a desk-judgment audit of the full $57$-problem LCB-medium pool,
classifying each problem by how hard it is to solve \emph{without any
explicit loop} (the loose rule):

\begin{center}
\begin{tabular}{lr}
\toprule
tier                                                 & count \\
\midrule
closed-form (no iteration needed)                     & $6$  \\
builtin-reducible (one-line \texttt{sum}/\texttt{map}/comprehension solves it) & $48$ \\
genuinely irreducible (real algorithmic iteration unavoidable)             & $\mathbf{3}$ \\
\bottomrule
\end{tabular}
\end{center}

\noindent
The three irreducibles are a breadth-first search for a minimum-edge
cycle, a multi-source breadth-first search on a grid with walls, and an
interval dynamic-program with a non-local recurrence. Every other problem
in the pool admits an answer that obeys the loose rule but pushes
iteration into a built-in. So a model that learns to write
\texttt{sum(map(int, input().split()))} instead of a \texttt{for} loop
has not internalised an aversion to iteration --- it has internalised the
\emph{shape} of one. The compliance signal we are training on is partly
real and partly a re-labelling. Worse, the constraint we are pretending
to enforce is mostly not biting on this benchmark.

\subsection*{Horn B --- the strict rule is harder than what we told the model, and ambiguous in places}

The strict judge fixes the toothlessness but breaks two other things.
First, it enforces a stricter rule than the one in the prompt. The model
was told ``no \texttt{for}, no \texttt{while}, no recursion''; the strict
judge additionally penalises ``and no \texttt{map}, no \texttt{reduce},
no \texttt{itertools}, no \texttt{sum} over a generator, no
\texttt{next}.'' A failure to satisfy that broader rule may be the
model's genuine inability to comply, or it may be the model honestly
following the rule it was actually given. We cannot tell these apart
from the output alone.

Second, the strict judge has to make several judgment calls that an
independent annotator might draw slightly differently:

\begin{itemize}
\item Is \texttt{sum([1,2,3])} iteration? (A fold over three elements ---
      yes by any honest reading.)
\item Is \texttt{sum(range(n))} iteration? (Also equal to $n(n{-}1)/2$.)
\item Is \texttt{min(s)} iteration? (It must visit every element.)
\item Is \texttt{next(iter(s))} iteration? (Once.)
\item Is \texttt{functools.reduce(op, xs, init)} recursion? (Operationally
      a left fold; semantically a catamorphism on a list, which is
      indistinguishable from a tail-recursive accumulator.)
\end{itemize}

The current strict judge takes positions on all of these. The positions
are defensible but not canonical, and a second person with the same
intent would draw the lines a little differently. To the extent
compliance under strict depends on these calls, the experiment starts
measuring the judge as much as the model.

\subsection*{The shape of the conflict}

\begin{center}
\begin{tabular}{lll}
\toprule
                       & \textbf{loose}              & \textbf{strict} \\
\midrule
where the boundary is  & abstract syntax tree       & multiple judgment calls \\
reproducible by anyone & yes                        & annotator-dependent     \\
how much it bites      & $3$ of $57$ problems       & nearly all of them      \\
matches the prompt     & yes                        & no (broader)            \\
risk                   & vacuous compliance         & measures the judge      \\
\bottomrule
\end{tabular}
\end{center}

\noindent
The current code reports both verdicts side-by-side and explicitly flags
results whose statistical significance depends on the choice. That is the
right thing to do \emph{methodologically} --- we do not have to commit
before looking. But it is not a \emph{scientific} resolution. Any headline
statement of the form ``training method X moved compliance from $a$ to $b$''
needs one of these two judges to be the primary one; the other becomes a
sensitivity check.

\section{What the next decisions are}
\label{sec:plan}

The follow-up quadrant experiment was specifically designed to put pressure
on the loose/strict question. The idea is to focus on a single
\emph{genuinely irreducible} problem --- one of the three in the audit --- and
ask: when R-SFT has trained the model to ``avoid loops,'' and the problem
makes loop-avoidance honestly impossible, does the model produce a confident
loose-compliant answer that is in fact strict-violating? In other words,
does the training induce a covering-up behaviour? And is that covering-up
visible in the model's internal activations?

This is a well-posed experiment because the loose/strict gap is forced ---
on an irreducible problem, the only way to be loose-compliant without being
strict-violating is to fail the tests. The contrast across the $12$ samples
per arm is what produces signal; the activation read is what makes it more
than a behavioural curiosity.

But the experiment cannot run cleanly until we decide three things:

\begin{enumerate}
\itemsep0pt
\item \textbf{Which judge is the headline number?} Loose is what we told the
      model and is reproducible; strict is what the rule actually means but
      involves judgment calls. A third option is to report both and treat
      the divergence as the finding --- but that weakens any single headline
      claim.
\item \textbf{Is the LCB-medium pool the right benchmark?} Three irreducible
      problems is below the $\sim\! 8$ threshold the project set itself for
      doing a re-cohort study. Switching to the hard slice (a one-line change
      in the data loader) should yield $\sim\! 50$ more irreducible problems
      and let us re-run Phases 1--2 on a benchmark where the constraint
      genuinely bites.
\item \textbf{Should the rule we tell the model match the rule we measure?}
      Under strict, the honest version of the experiment puts the broader
      rule in the prompt --- ``no \texttt{for}, no \texttt{while}, no
      recursion, no \texttt{map}, no \texttt{reduce}, no \texttt{itertools},
      no \texttt{sum}/\texttt{any}/\texttt{all}/\texttt{min}/\texttt{max}
      over a generator, no \texttt{next}.'' That is a different experiment
      from the one we ran for Phases 1--2, and the existing numbers do not
      transfer to it.
\end{enumerate}

\section{One-paragraph summary}

The training pipeline works, the rationale-SFT result is real, the next
experiment is queued and waiting. The blocker is conceptual rather than
technical. The rule that defines the whole project admits a permissive
reading under which roughly $95\%$ of the benchmark is trivially compliant
(and therefore the rule is barely studying anything), and a faithful
reading that is broader than what the model was told and unavoidably
involves judgment calls. The project needs to commit to one of these two
readings before the next launch --- or commit to the harder thing,
publishing both and treating the gap between them as the finding.

\end{document}
